% ===== 原文 pp.30–31（PDF pp.31–32）=====

% [原文 p.30]
FI 参数空间 $\R^2=\{(\zeta^1,\zeta^2)\}$ 被分成三个相，如图 3 所示。相 I 是第一象限，$\zeta^{1,2}>0$；相 II 是腔室 $\zeta^1<0$ 且 $\zeta^1+n\zeta^2>0$；相 III 是图中阴影所表示的其余区域。

在相 I 中，$(\phi_1,\phi_2)$ 与 $(\phi_3,\phi_4)$ 都必须非零，规范群被完全破缺。理论约化为靶空间为 Hirzebruch 曲面 $\mathbb{F}_n=\mathbb{P}(\cO\oplus\cO(-n))$ 的 sigma 模型。该曲面在 $n=0,1$ 时为 Fano，在 $n=2$ 时为 nef，在 $n\ge 3$ 时为非 Fano。非零的 Hodge 数为 $h^{0,0}=1$，$h^{1,1}=2$，$h^{2,2}=1$。

在相 II 中（该相只在 $n\ge 1$ 时存在），$(\phi_1,\phi_2,\phi_3)$ 与 $\phi_4$ 都必须非零。$\phi_4$ 的非零值把规范群破缺到子群 $\{(g,g^n)\}\cong U(1)$，在该子群下 $\phi_1,\phi_2,\phi_3$ 的荷分别为 $1,1,n$。除去 $\phi_1=\phi_2=0$ 的那些点（在那里它被破缺到 $\Z_n$ 子群）之外，这个剩余群被 $(\phi_1,\phi_2,\phi_3)$ 的非零值完全破缺。理论约化为靶空间为加权射影空间 $\WP^2_{1,1,n}$ 的 sigma 模型。当 $n\ge 2$ 时，空间 $\WP^2_{1,1,n}$ 有一个 $\C^2/\Z_n$ 型的轨形点，其中 $\Z_n$ 通过 $(z_1,z_2)\mapsto(\ee^{2\pi\ii/n}z_1,\ee^{2\pi\ii/n}z_2)$ 作用在 $\C^2$ 上。注意相 I 中的 $\mathbb{F}_n$ 可以看作这个轨形奇点的一个消解。事实上，$\mathbb{F}_n\to\WP^2_{1,1,n}$ 正是模型 $\mathrm{T}^{U(1)}_{N,d}$（取 $N=2$、$d=n$）中出现过的消解 $\mathrm{tot}[\cO_{\mathbb{P}^1}(-n)]\to\C^2/\Z_n$ 的紧版本。

在相 III 中，方程 (4.34) 没有解。也就是说，超对称在经典层面就被破缺了。

在量子理论中，FI 参数按下式跑动：$\zeta^{1'}=\zeta^1+(2-n)\log(\mu'/\mu)$，$\zeta^{2'}=\zeta^2+2\log(\mu'/\mu)$。对任何 $n$，$\zeta^2$ 都从大的正值跑到大的负值；而 $\zeta^1$ 的跑动方向取决于 $n\le 1$（从正到负）、$n=2$（不跑动）、$n\ge 3$（从负到正）。这三种情形下面分别讨论。在大 $\sigma$ 值处计算有效扭超势，留给读者作练习。这里我们只写下那里的真空方程：
\begin{equation}
    \sigma_1^2=q_1\mu^{2-n}(-n\sigma_1+\sigma_2)^n,\qquad
    \sigma_2(-n\sigma_1+\sigma_2)=q_2\mu^2, \tag{4.35}
\end{equation}
其中 $q_a:=\ee^{-t^a}$。$n=0,1$ 时有四个解；$n=2$ 时一般有四个解；$n\ge 3$ 时有 $n+2$ 个解（见下文）。
% [Tutor Note: p.30 式 (4.35) 可用来自我核对：取 n=0 时它化为 \sigma_1^2=q_1\mu^2 与 \sigma_2^2=q_2\mu^2，
% 恰好给出 2\times2=4 个解，与下文"n=0,1 时有四个解"一致；取 n=1 时第二式不变、第一式化为
% \sigma_1^2=q_1\mu(-\sigma_1+\sigma_2)。另注意 q_a:=\ee^{-t^a} 中的 t^a 是复化 FI-theta 参数
% （t^a=\zeta^a-\ii\theta^a，与印刷 p.23 的 $t=\zeta-\ii\theta$ 同一约定），这一点在印刷 p.34 的 D-项真空方程 \sum_i Q_i^a|\phi_i|^2=\zeta^a 处得到印证。
% 原文"（see blow）"是"see below"之误。]

\begin{figure}[htbp]
\centering
\includegraphics[width=0.9\textwidth]{fig4.png}
\figtitle{4}{$\mathrm{H}_0$（左）与 $\mathrm{H}_1$（右）的 RG 流。}
\end{figure}

\para{n = 0, 1} FI 参数的流动如图 4 所示。在两种情形中，任何流的紫外极限都在相 I 中，而相 I 分别对应于 $\mathbb{F}_0=\CP^1\times\CP^1$（$n=0$）或 $\mathbb{F}_1$（即 $\CP^2$ 的一点爆破，$n=1$）。因此，GLSM 可以等同于以这些空间为靶空间的 sigma 模型。注意这两个靶空间都是 Fano 流形，因此相应的 sigma 模型是紫外完备的。在模型 $\mathrm{H}_1$ 中，有些流在中间能标处会穿过相 II。回忆一下，此情形中相 II 对应于 $\WP^2_{1,1,1}=\CP^2$。我们来讨论能否把 $\CP^2$ sigma 模型认作模型 $\mathrm{H}_1$ 的%
% [Tutor Note: 原文 p.30 的 "one pont blow up" 是 "one point blow up"（一点爆破）之误；\mathbb{F}_1
% 正是一点爆破的 \CP^2，其 h^{1,1}=2 与上文"非零 Hodge 数 h^{1,1}=2"相符。另外原文此处写作
% WCP^2_{1,1,1}（W+\CP，多出一个 C），与本章其余各处 \WP^2_{1,1,n} 的排印不一致，本译按统一记号
% \WP 处理；核对方法：p.30 该处图像中 W 与 P 之间确有一个双线 C。]
% [原文 p.31]
一部分。我们写出 FI-theta 参数为 $q_1=\Lambda/\mu$、$q_2=q(\Lambda/\mu)^2$，其中 $q$ 与 $\Lambda$ 是该理论的重整化群不变量。注意当 $q_1\gg 1$ 且 $q_1q_2\ll 1$ 时理论处于相 II，也就是处在能标窗口 $q^{1/3}\Lambda\ll\mu\ll\Lambda$ 内，而这只有在 $q\ll 1$ 时才可能。若取极限 $q\to 0$ 并固定 $\Lambda^3_{\mathbb{P}^2}:=q\Lambda^3$，则窗口变为 $\Lambda_{\mathbb{P}^2}\ll\mu<\infty$，与 $\CP^2$ sigma 模型中的情形一样。方程 (4.35) 可以写成
\begin{equation}
    \sigma_2=\sigma_1+\sigma_1^2/\Lambda,\qquad
    \sigma_1^3(\sigma_1+\Lambda)=q\Lambda^4. \tag{4.36}
\end{equation}
% [Tutor Note: 式 (4.36) 可由 (4.35) 代入 n=1、q_1=\Lambda/\mu、q_2=q(\Lambda/\mu)^2 直接验证：
% (4.35) 第一式给 \sigma_1^2=\Lambda(-\sigma_1+\sigma_2)，即 \sigma_2=\sigma_1+\sigma_1^2/\Lambda；
% 再把 \sigma_2=\sigma_1(\sigma_1+\Lambda)/\Lambda 与 -\sigma_1+\sigma_2=\sigma_1^2/\Lambda 代入第二式
% 得 \sigma_1^3(\sigma_1+\Lambda)=q\Lambda^4。这条推导同时说明下文 \sigma_1\sim-\Lambda 的单个解为何
% 在 q\to0 时跑到无穷远（该支上 \sigma_1+\Lambda 很小）。]
当 $q\ll 1$ 时，有三个解满足 $\sigma_1^3\sim q\Lambda^3$、$\sigma_1\sim\sigma_2$，另有一个解满足 $\sigma_1\sim-\Lambda$、$\sigma_2\sim q\Lambda$。在上述标度极限下，那三个解保持有限（$\sigma_1^3\sim\Lambda^3_{\mathbb{P}^2}$），而单个解则跑到无穷远。这些表明，该标度极限把 $\CP^2$ sigma 模型从 $\mathrm{H}_1$ 中取了出来。

\para{n = 2} FI 参数的流动如图 5 左图所示。$q_1$ 是重整化群不变量，但 $q_2$ 按 $q_2=(\Lambda/\mu)^2$ 跑动。当 $q_1\ll 1$ 时理论描述 $\mathbb{F}_2$ sigma 模型，当 $q_1\gg 1$ 时描述 $\WP^2_{1,1,2}$ sigma 模型。对 $q_1\ne 1/4$，真空方程有四个解
\begin{equation}
    \sigma_1=\frac{\epsilon_2\epsilon_1 q_1^{\frac12}\Lambda}{(1+2\epsilon_1 q_1^{\frac12})^{\frac12}},\qquad
    \sigma_2=\epsilon_2\Lambda(1+2\epsilon_1 q_1^{\frac12})^{\frac12},\qquad
    \epsilon_1,\epsilon_2\in\{\pm1\}. \tag{4.37}
\end{equation}
% [Tutor Note: 用 (4.37) 可以核对下一句"两个解跑到无穷远"：q_1\to1/4 时分母
% (1+2\epsilon_1q_1^{1/2})^{1/2}\to(1+\epsilon_1)^{1/2}，唯一发散的是 \epsilon_1=-1 的两支
% （\epsilon_2=\pm1），共两个解；\epsilon_1=+1 的两支保持有限。这也说明四解在 q_1=1/4 处两两
% 分歧，故 q_1=1/4 确为判别轨迹。]
当 $q_1\to 1/4$ 时，其中两个解跑到无穷远；$q_1=1/4$ 可以看作一条判别轨迹。$\WP^2_{1,1,2}$ 与 $\mathbb{F}_2$ 两个模型沿 $q_1$ 空间中一条从 $q_1\gg 1$ 到 $q_1\ll 1$、避开 $q_1=1/4$ 的路径相连。在此过程中，$\WP^2_{1,1,2}$ 的 $\C^2/\Z_2$ 奇点（$A_1$ 型奇点）通过插入一个 $\CP^1$ 而被消解。我们看到，McKay 对应被实现为该紧模型的一部分。

\para{n $\ge$ 3} FI 参数的流动如图 5 右图所示。任何流的紫外极限都在相 II 中。因此该 GLSM 可以等同于 $\WP^2_{1,1,n}$ sigma 模型。
